\onecolumn

\begin{table*}[!t]
\centering
\section{Tabla de artículos}
\renewcommand{\arraystretch}{1.3}

\resizebox{\textwidth}{!}{
\begin{tabular}{|p{3cm}|p{4cm}|p{1.5cm}|p{3cm}|p{6cm}|}
\hline
\textbf{Título} & \textbf{Autores} & \textbf{Año} & \textbf{Revista} & \textbf{Resumen} \\ 
\hline

1. Estimated association between dwelling soil contamination and internal radiation contamination levels after the 2011 Fukushima Daiichi nuclear accident in Japan
& M. Tsubokura; S. Nomura; K. Sakaihara; S. Kato; C. Leppold; T. Furutani; T. Morita; T. Oikawa; Y. Kanazawa
& 2016
& BMJ Journals
& Measurement of soil contamination levels has been considered a feasible method for dose estimation of internal radiation exposure following the Chernobyl disaster by means of aggregate transfer factors; however, it is still unclear whether the estimation of internal contamination based on soil contamination levels is universally valid or incident specific. Low levels of internal and soil contamination were not associated, and only loose/small associations were observed in areas with slightly higher levels of soil contamination in Fukushima, representing a clear difference from the strong associations found in post-disaster Chernobyl.\\
\hline

2. Measurement of radioactivity in building materials – problems caused by possible disequilibrium in natural decay series
& B. Michalik; G. de With; Schroeyers
& 2018
& Construction and Building Materials
& Current regulations in the European Union require the calculation of an activity concentration index (index I) using the activity concentrations of $^{226}$Ra, $^{232}$Th, and $^{40}$K. Not all of these radionuclides are directly measurable by gamma spectrometry, so additional assumptions about secular equilibrium in the uranium and thorium decay series are required. These assumptions are often invalid in Naturally Occurring Radioactive Materials (NORM), where long-term disequilibrium is commonly observed. This can lead to either underestimation or overestimation of the index. \\
\hline

3. Temporal Change in Radiological Environments on Land after the Fukushima Daiichi Nuclear Power Plant Accident
& K. Saito; S. Mikami; M. Andoh; N. Matsuda; S. Kinase; S. Tsuda; T. Sato; A. Seki; Y. Sanada; H. Wainwright-Murakami; K, Yoshimura; H. Takemiya; J. Takahashi; H. Kato; Y. Onda
& 2019
& Journal of Radiation Protection and Research
& Massive environmental monitoring has been conducted continuously since the Fukushima Dai-ichi Nuclear Power accident in March of 2011 by different monitoring methods that have different features together with migration studies of radiocesium in diverse environments. These results have clarified the characteristics of radiological environments and their temporal change around the Fukushima site. Radiation levels around the Fukushima site have decreased greatly over time. The decreasing trend was found to change variously according to local conditions. \\
\hline

4. Numerical simulation of a temporary repository of radioactive material
& E. de la Cruz-Sánchez; J. Klapp; E. Mayoral-Villa; R. González-Galán; A. M. Gómez-Torres
& 2017
& International Journal of Nuclear Energy Science and Technology
& This work presents a bidimensional numerical simulation model to study the dispersion of contaminants through saturated porous media using the finite element method (FEM). The model is applied to analyze the transport of radioisotopes in a temporary nuclear repository located in the Vadose Zone at Peña Blanca, Mexico. The main result is the development of a validated model describing the migration of radionuclides in saturated porous media with fractures.\\
\hline

5. Environmental behaviour of radioactive particles from Chernobyl
& B. Salbub; S. Levchuk; V. Protsak; I. Maloshtan; C. Simonucci; C. Courbet; H.L. Nguyen; N. Sanzharova; V. Zabrotsky
& 2019
& Journal of Environmental Radioactivity
& This study analyzes the long-term environmental behaviour of radioactive particles released during the Chernobyl accident. The particles were deposited in sandy topsoil in the Ivankiv district of Kyiv region (Ukraine), in trench waste materials from the Red Forest, and in bottom sediments of the cooling pond. Results show that partial dissolution of fuel particles with low chemical stability (UO$_{2+x}$) has occurred, along with partial dissolution of more stable UO$_2$ particles. This indicates progressive radiological stabilization of the environment. The mobile fraction of radionuclides is expected to decrease over time, while the biological availability of $^{90}$Sr in topsoil has reached a maximum and is expected to decline.\\
\hline

6.	Fukushima and Chernobyl: Similarities and Differences of Radiocesium Behavior in the Soil–Water Environment
& A. Konoplev
& 2022
& Toxics
& In the wake of Chernobyl and Fukushima accidents, radiocesium has become a radionuclide of most environmental concern. The ease with which this radionuclide moves through the environment and is taken up by plants and animals is governed by its chemical forms and site-specific environmental characteristics. Disintegration of these particles in the environment is much slower than that of Chernobyl-derived fuel particles. The higher annual precipitation and steep slopes in Fukushima-contaminated areas are conducive to higher erosion and higher total radiocesium wash-off. \\
\hline

\end{tabular}
}
\end{table*}

\begin{table*}[!t]
\centering
\renewcommand{\arraystretch}{1.3}

\resizebox{\textwidth}{!}{
\begin{tabular}{|p{3cm}|p{4cm}|p{1.5cm}|p{3cm}|p{6cm}|}
\hline
\textbf{Título} & \textbf{Autores} & \textbf{Año} & \textbf{Revista} & \textbf{Resumen} \\ 
\hline

7.	Long-term study (1987–2023) on the distribution of 137Cs in soil following the Chernobyl nuclear accident: a comparison of temporal migration measurements and compartment model predictions
& I. Kaissas; A. Clouvas; M. Postatziis; S. Xanthos; M. Omirou
& 2023
& Radiation Protection Dosimetry
& After the Chernobyl accident, a designated area of $\sim 1000 m^2$ within the University farm of Aristotle University of Thessaloniki in Northern Greece was utilized as a test ground for radioecological measurements. The profile of 137Cs in the soil was monitored from 1987 to 2023, with soil samples collected in 5-cm-thick slices (layers) down to a depth of 30 cm. A compartment model was utilized to forecast the temporal evolution of fractional contributions of the different soil layers to the total deposition of $^{137}$Cs (0–30 cm). In this model, each soil layer is represented as a separate compartment. However, the use of a second compartment model with increasing transfer rates between consecutive soil layers did not align with the observed outcomes. This indicates that diffusion may not be the primary migration mechanism over the 36-y period covered by our study.\\
\hline

*Generational Garbage Collection and the Radioactive Decay Model
& W.D. Clinger; L.T. Hansen
& 1997
& PLDI Conference
& This work analyzes generational garbage collection using a radioactive decay model, where object lifetimes are assumed to follow an exponential probability distribution. Under this assumption, the age of an object provides no information about its future lifetime. The model shows that there is no rational basis for predicting which objects will live longer than others, but there is still a basis for deciding how many objects to promote and when to perform garbage collection. The analysis leads to a generational garbage collection strategy that does not rely on heuristics about object lifespan, providing insight into memory management for systems with unpredictable object lifetimes.\\
\hline

*Isotope Shifts, Radioactive Decay, and the Nuclear Forces
& W. Raven
&
& 2024
& This chapter explores the nuclear structure, focusing on how the number of neutrons affects atomic transition frequencies and nuclear stability. It discusses the role of nuclear forces in governing stability and radioactive decay. Various decay modes are analyzed, including neutron emission, proton emission, $\alpha$ decay, spontaneous fission, $\beta^-$ decay, $\beta^+$ decay, and electron capture, with emphasis on the conditions under which each process occurs. The nuclear shell model is also introduced to explain the energetics and stability of isotopes.\\
\hline

*The discovery of radioactivity
& P. Radvanyi; J. Villain
& 2017
& Comptes Rendus Physique
& This work describes the historical development of the discovery of radioactivity. Early studies of the penetrating power of radiation and its behavior under electric and magnetic fields revealed the complexity of nuclear radiation, consisting of three components: $\alpha$, $\beta$, and $\gamma$ radiation. The work highlights how difficult it was to understand the nature of radioactivity, which was later clarified by Ernest Rutherford and Frederick Soddy.\\
\hline

Global risk of radioactive fallout after major nuclear reactor accidents
& J. Leleiveld; D. Kundel; M.G. Lawrence
& 2011
& Atmospheric Chemistry and Physics
& Major nuclear reactor accidents are rare, yet their consequences are severe. This work examines what is meant by “rare” in this context and discusses lessons from the Chernobyl and Fukushima accidents. The study assesses the global risk of exposure to radioactive fallout due to atmospheric dispersion of gases and particles following severe nuclear accidents classified as International Nuclear Event Scale (INES) level 7. The analysis uses $^{137}$Cs and $^{131}$I as tracers of radioactive dispersion.\\
\hline

Radioactivity release from the Fukushima accident and its consequences: A review
& Y.H. Koo; Y.S. Yang; K.W. Song
& 2014
& Journal of Environmental Radioactivity
& The Fukushima accident in March 2011, triggered by a massive earthquake and tsunami, led to hydrogen explosions, core meltdowns, and the release of large amounts of radioactivity into the atmosphere and the Pacific Ocean. This review examines the accident from the perspective of radioactive release and environmental dispersion, as well as its impacts on public health, the economy, energy policy, international relations, and light-water reactor (LWR) fuel development.\\
\hline

\end{tabular}
}
\end{table*}

\begin{table*}[!t]
\centering
\renewcommand{\arraystretch}{1.3}

\resizebox{\textwidth}{!}{
\begin{tabular}{|p{3cm}|p{4cm}|p{1.5cm}|p{3cm}|p{6cm}|}
\hline
\textbf{Título} & \textbf{Autores} & \textbf{Año} & \textbf{Revista} & \textbf{Resumen} \\ 
\hline

Evidence of the radioactive fallout in France due to the Fukushima nuclear accident
& O. Evrard; P. Van Beek; D. Gateuille; V. Pont; I. Lefèvre; B. Lansard; P. Bonté
& 2012
& Journal of Environmental Radioactivity
& Radioactive fallout from the Fukushima nuclear accident in Japan was detected in environmental samples collected in France, demonstrating long-range atmospheric transport of radionuclides following the reactor accident. \\
\hline

*Spectral content of $^{22}$Na / $^{44}$Ti decay data
& D. O'Keefe; B. Morreale; P.A. Sturrock; et al.
& 2012
& Astrophysics and Space Science
& This work analyzes the spectral content associated with the radioactive decay of $^{22}$Na and $^{44}$Ti, focusing on gamma-ray emission lines relevant for astrophysical observations. The study investigates how decay signatures can be used to extract physical information from spectral data and improve the interpretation of nuclear decay processes in observational contexts. \\
\hline

*Radioactive Decay 
& Paola Cappellaro 
& 2023 
& Physics LibreTexts (MIT OpenCourseWare, educational resource) 
& Describes radioactive decay as the process by which an unstable nucleus spontaneously loses energy by emitting particles and ionizing radiation; introduces the concept of nuclear instability and its spontaneous nature within the framework of basic nuclear physics. \\
\hline

*Isotope Shifts, Radioactive Decay, and the Nuclear Forces
& W. Raven
& 2024
& (Book chapter / not specified)
& This chapter explores nuclear structure, focusing on how the number of neutrons affects atomic transition frequencies and nuclear stability. It discusses the role of nuclear forces in governing stability and radioactive decay. Various decay modes are analyzed, including neutron emission, proton emission, $\alpha$ decay, spontaneous fission, $\beta^-$ decay, $\beta^+$ decay, and electron capture.\\
\hline


\end{tabular}
}
\end{table*}


